\renewcommand \arraystretch{1.05}
\begin{table*}[t]
    \setlength{\tabcolsep}{4.5pt}
    \small
    \centering
    \vspace{5pt}
    \begin{tabular}{lccccccccccc}
        \toprule
          \multirow{2}{*}{\textbf{PPL $\downarrow$}}   & \multirow{2}{*}{FP16} & \multirow{2}{*}{RTN}   & \multicolumn{3}{c}{FP16\% (based on act.)}  & \multicolumn{3}{c}{FP16\% (based on W)} &  \multicolumn{3}{c}{FP16\% (random)} \\ \cmidrule(lr){4-6} \cmidrule(lr){7-9}  \cmidrule(lr){10-12}
         & & (w3-g128) & 0.1\% & 1\% & 3\% &   0.1\% & 1\% & 3\% &  0.1\% & 1\% & 3\%   \\
        \midrule
        OPT-1.3B  & 14.62 & 119.00 & 25.03 & \textcolor{codegreen}{16.91} & \textcolor{codegreen}{16.68} & 108.71 & 98.55 & 98.08 & 119.76 & 109.38 & 61.49 \\
        OPT-6.7B & 10.86 & 23.54 & \textcolor{codegreen}{11.58} & \textcolor{codegreen}{11.39} & \textcolor{codegreen}{11.36} & 23.41 & 22.37 & 22.45 & 23.54 & 24.23 & 24.22 \\
        OPT-13B &  10.13 & 46.04 &\textcolor{codegreen}{10.51} & \textcolor{codegreen}{10.43} & \textcolor{codegreen}{10.42} & 46.07 & 48.96 & 54.49 & 44.87 & 42.00 & 39.71 \\
        \bottomrule
    \end{tabular}
    \caption{Keeping a small fraction of weights (0.1\%-1\%) in FP16 significantly improves the performance of the quantized models over round-to-nearest (RTN). It is only effective when we select the important weights in FP16 by looking at \emph{activation} distribution instead of \emph{weight} distribution. We highlight results with a decent perplexity in \textcolor{codegreen}{green}.  
    We used INT3 quantization with a group size of 128 and measured the WikiText perplexity ($\downarrow$).  
    } 
    \label{tab:fp_ratio}
\end{table*}
